\documentclass[12pt]{article}
\usepackage[letterpaper, margin=1in]{geometry}
\usepackage{titlesec}
\usepackage{titling}
\usepackage{graphicx}
\usepackage{url}
\usepackage{amsmath}
\usepackage{color}
\usepackage[colorlinks=true, urlcolor=blue, linkcolor=blue, citecolor=blue]{hyperref}
\bibliographystyle{apalike}
\usepackage{biblatex}


% Title formatting courtesy of Damein Koon
\pretitle{\LARGE\bfseries}
\posttitle{\par\vspace{2em}}
\preauthor{\large}
\postauthor{\par\vspace{1em}}
\predate{\large}
\postdate{\par\vspace{3em}}

% Corrected title hooks for vertical and horizontal centering
\renewcommand\maketitlehooka{\null\vfill\begin{center}}
\renewcommand\maketitlehookd{\end{center}\vfill\null}

\title{Quasi-Normal Modes (QNMs) of Black Holes Project Proposal}
\author{Krupa Pothiwala \\[0.5em]
March 2025}
\date{}


\begin{document}
\maketitle

\newpage
\section{Goal of the Project}
\label{Goal}
When a black hole merger happens, during the ringdown phase the black hole isn't a clean Kerr black hole. It is a perturbed, distorted object that radiates gravitational waves as it settles down. The waves it is emitting are called Quasi-Normal modes. Quasi because they're not truly normal modes, as energy leaks through, so they are damped. For a Kerr black hole these modes only depends on its mass $M$ and spin $a$. In this project, we will compute these frequencies. 
\\
\\
The physics that I will need in order to work on this project is: Black Hole Perturbation Theory, The Regge-Wheeler \& Zerilli Equations, and The Teukolsky Equation.

\section{Computational Part}
\label{Computer}
On the computational side, I will implement Leaver's continued fraction method in Julia to numerically solve for these complex QNM frequencies, validating my results against published tables and ultimately plotting how the spectrum changes with black hole spin.

\section{Literature Part}
\label{Lit}
I may not use all of these directly, but I will definitely need to use them to gain background knowledge.
\begin{enumerate}
    \item \cite{ReggeWheeler1957} - First perturbation theory of Schwarzschild
    \item \cite{Zerilli1970} - Polar perturbations, completes Schwarzschild case
    \item \cite{Teukolsky1972} \& \cite{Teukolsky1973} - The master equation for Kerr
    \item \cite{PressTeukolsky1973} - Floating orbits, physical interpretation
    \item \cite{Vishveshwara1970} - First numerical evidence of QNMs
    \item \cite{Leaver1985} - Continued fraction method
    \item \cite{KokkotasSchmidt1999} - The canonical QNM review, very readable
    \item \cite{BertiCardosoStarinets2009} - Comprehensive modern review; has QNM tables to validate against
    \item \cite{Nollert1999} - Good physical intuition on QNMs
    \item \cite{GW150914} - First detection, ringdown section
    \item \cite{Isi2019} - First overtone detection in GW150914 ringdown
    \item \cite{CotaescuGazeau} - Spectroscopy and no-hair tests
\end{enumerate}


\end{document}